
\subsubsection{Interpretation of Key Findings}

\textbf{Perfect Scale Invariance.} The coefficient of variation = 0.00 across ResNet-{18,34,50,101,152} indicates parameter-mass ratio R is perfectly invariant within homogeneous architecture families. ResNet scales by stacking more residual blocks with identical convolutional structure, preserving the ratio of convolutional to linear parameters regardless of depth. This suggests R captures architectural computation paradigm rather than incidental properties that vary with model size.

\textbf{Exceptional Separation.} Cohen's d = 3.202 between CNN and Transformer R distributions is extraordinarily large---conventional thresholds define d > 0.8 as large effect. This magnitude indicates convolutional and attention-based paradigms impose fundamentally distinct parameter allocation strategies with near-zero overlap. These are not continuous variations---they are discrete architectural paradigms separable with simple ratio.

\textbf{Normalization Fingerprinting.} The 0\% CNN violation rate for BatchNorm usage indicates not a single CNN model used LayerNorm as dominant normalization. This is mechanistic coupling: BatchNorm enforces spatial normalization suited for image data where spatial locations share statistics, while LayerNorm enforces token-wise normalization suited for sequential data where dimensions vary independently. Empirical validation confirms Chun (2026) theoretical prediction without requiring forward passes.

\textbf{Naming Alignment Failure.} The paradoxical result that TIMM naming alignment was only 40\% yet classification succeeded (88.89\% accuracy) proves that features extract structural information from checkpoint tensors independent of naming. This robustness to labeling noise is valuable for model zoo management where naming may be inconsistent.

\subsubsection{Comparison to Related Work}

\textbf{Kofinas et al. (2024).} This approach achieves comparable architectural understanding (88.89\% family classification) to graph neural network methods but with $100\times$ faster extraction and full interpretability. Kofinas requires constructing computational graphs, training GNNs with message passing, and 50+ hours implementation. This checkpoint inspection requires loading dictionary, counting layer names, and querying tensor shapes---total implementation <6 hours. The tradeoff: GNNs can potentially learn more complex weight-space representations for tasks beyond family classification.

\textbf{Operationalizing Theory.} Prior work provided theoretical foundations (Chun: LayerNorm vs BatchNorm geometry) and empirical observations (Fang: diverged parameter scales) but did not demonstrate practical applications. This work operationalizes these insights: normalization layer counts exploit Chun's predictions without requiring forward passes to measure geometric effects, while parameter-mass ratio exploits Fang's observation by computing explicit allocation ratio.

\textbf{Zhang \& Abdulla (2023).} Zhang extracted architectural information from BatchNorm kernel weights but required forward passes to populate running statistics. This checkpoint-only constraint demands features computable from state\_dict alone, leading to counting normalization layer types rather than analyzing parameters. The result is faster extraction (no model instantiation) at cost of lower-resolution information. For family classification, counts suffice.

\subsubsection{Limitations}

\textbf{L1: NormFree Network Failure.} NormFree architectures (NFNet, NormFree-ResNet) achieve 0\% classification accuracy because they replace normalization layers with scaled weight standardization. With bn\_count = ln\_count = gn\_count = 0, the method falls back to parameter-mass ratio alone, which is insufficient for these paradigms. This limitation is precisely characterized: NormFree networks occupy a third paradigm the binary fingerprinting scheme cannot capture. Mitigation would extend features with weight distribution statistics or activation layer counts. NormFree architectures are rare in production (VGG-16 historical, NFNets niche). Method successfully handles many other edge cases (SENet, RegNet, ViT-Extreme all 100\% accurate).

\textbf{L2: Small Validation Set.} Validation set contains only 18 models, limiting detection of rare failure modes and reducing precision of accuracy estimates. Primary metrics exceed thresholds with comfortable margins (+8.89pp for P1), providing robust directional evidence despite wide intervals. Mitigation would expand to full TIMM zoo validation (1000+ models) using checkpoint extraction feasibility (1.05s per model, 17 minutes total). Proof-of-concept prioritized implementation speed (<8 hours) over exhaustive validation.

\textbf{L3: Transformer Scale Invariance Unverified in H-M2.} While CNN scale invariance is directly confirmed (ResNet CV=0.00), Transformer scale invariance is supported by training set analysis but not independently reconfirmed in h-m2 validation due to insufficient scale-family models in validation set (only 1 ViT model). This is validation oversight, not conceptual gap---Transformer architecture modularity mechanistically suggests scale invariance similar to ResNet.

\textbf{L4: Scope Limited to Vision Models.} All experiments use vision models from TIMM zoo. Generalization to language models, audio models, or non-TIMM architectures is unverified. Language models predominantly use LayerNorm regardless of whether they are CNN-like or Transformer-like, potentially violating normalization fingerprinting assumption. Extension to language models would require alternative features.

\subsubsection{Broader Impact}

\textbf{Model Zoo Management.} Zero GPU requirement and 1.05s per-model extraction enable architecture verification on commodity hardware. Practitioners managing large model repositories can automatically classify architectures in minutes, facilitating architecture-aware model selection for meta-learning, automated filtering for model merging, and architecture-based organization independent of naming conventions.

\textbf{Paradigm Shift.} Traditional model analysis requires instantiation and forward passes. This checkpoint-only approach enables understanding models from weights alone, opening possibilities for checkpoint archaeology---extracting information about training dynamics, dataset biases, or model provenance without execution. This shift reduces barriers when execution is expensive or risky.

\textbf{Mechanistic Feature Engineering.} Success with hand-crafted features guided by mechanistic understanding demonstrates alternative to end-to-end neural approaches. When domain knowledge provides mechanistic explanations for discriminative signals, explicit feature engineering can match neural methods with better interpretability and efficiency.
